\section*{Bab \ref{ch:b3:surfaces-catalysis} --- Permukaan, Adsorpsi, dan Katalisis Heterogen}

\begin{solution}{exo:b3:surfaces-catalysis:1}
$K = \theta/((1 - \theta)p) = 0.40/(0.60 \times 2.0) = \qty{0.33}{kPa^{-1}}$; pada \qty{10}{kPa}, $\theta = 3.33/4.33 = 0.77$.
\end{solution}

\begin{solution}{exo:b3:surfaces-catalysis:2}
$-20$: \omterm{def:b3:surfaces-catalysis:physi-chemi}{fisisorpsi} (berorde entalpi kondensasi); $-150$: \omterm{def:b3:surfaces-catalysis:physi-chemi}{kemisorpsi}, yang
merupakan satu-satunya yang dapat bersifat disosiatif.
\end{solution}

\begin{solution}{exo:b3:surfaces-catalysis:3}
(100): satu situs puncak, dua situs jembatan (empat ikatan yang masing-masing
dimiliki bersama oleh dua atom), satu rongga lipat empat (empat rongga yang
masing-masing dimiliki bersama oleh empat atom). (111): satu situs puncak, tiga
situs jembatan, dua rongga lipat tiga (enam segitiga di sekeliling setiap
atom, yang masing-masing dimiliki bersama oleh tiga atom).
\end{solution}

\begin{solution}{exo:b3:surfaces-catalysis:4}
$\num{3.0e-6}/\num{1.5e-5} = \qty{0.20}{s^{-1}}$.
\end{solution}

\begin{solution}{exo:b3:surfaces-catalysis:5}
$\theta/(1 - \theta) = (Kp)^{1/2}$: $0.111$ pada \qty{1.0}{kPa}, menjadi dua kali lipat pada \qty{4.0}{kPa}: $0.222$, sehingga $\theta = 0.18$.
\end{solution}

\begin{solution}{exo:b3:surfaces-catalysis:6}
Penyebut $1 + 10 + 2.5 = 13.5$: $\theta_{\ce{CO}} = 0.74$, $\theta_{\ce{O2}} = 0.19$, situs bebas 0.07.
Laju Langmuir--Hinshelwood, yang sebanding dengan $\theta_{\ce{CO}}\theta_{\ce{O2}}$, dibatasi oleh
langkanya oksigen pada permukaan yang dipenuhi CO.
\end{solution}

\begin{solution}{exo:b3:surfaces-catalysis:7}
$\Delta_{\mathrm{ads}}H = -R\ln(8.0/1.0)/(1/300 - 1/350) = -8.314 \times 2.079/\num{4.76e-4} = \qty{-36}{kJ/mol}$.
\end{solution}

\begin{solution}{exo:b3:surfaces-catalysis:8}
$n_{\mathrm m} = 120/22\,414 = \qty{5.35e-3}{mol}$; luas $= \num{5.35e-3} \times \num{6.022e23} \times \num{0.162e-18} =
\qty{522}{m^2.g^{-1}}$.
\end{solution}

\begin{solution}{exo:b3:surfaces-catalysis:9}
Langmuir--Hinshelwood: kedua pereaksi harus teradsorpsi pada situs-situs yang
bertetangga; CO terikat kuat dan, pada tekanan tinggi, tidak menyisakan tempat
bagi oksigen.
\end{solution}

\begin{solution}{exo:b3:surfaces-catalysis:10}
Tutupan rendah: $100 - 60 = \qty{40}{kJ/mol}$; tutupan tinggi: \qty{100}{kJ/mol}. Plot $\ln r$
terhadap $1/T$ curam (kemiringan $-100/R$) pada suhu rendah, ketika permukaannya
penuh, dan lebih landai (kemiringan $-40/R$) pada suhu tinggi, ketika permukaannya
kosong: sebuah kurva yang melengkung di antara keduanya.
\end{solution}

\begin{solution}{exo:b3:surfaces-catalysis:11}
$4 \times 0.10 = 0.40$ bagian situs tersumbat: 60\,\% aktivitas tersisa untuk satu
situs, dan sekitar $0.60^2 = 36\,\%$ untuk dua situs bebas yang bertetangga.
\end{solution}

\begin{solution}{exo:b3:surfaces-catalysis:12}
Logam yang tengah (\qty{-250}{kJ/mol}): logam pertama mengikat nitrogen terlalu lemah untuk memutuskan \ce{N2}
dengan mudah, sedangkan logam ketiga terlalu kuat, sehingga atom-atom nitrogen
tetap tinggal di permukaan dan menyumbatnya.
\end{solution}

\begin{solution}{pb:b3:surfaces-catalysis:1}
\textbf{1.} Nitrogen mengembun di dekat \qty{77}{K} pada tekanan atmosfer: seluruh rentang
tekanan relatif dapat dicapai, dan multilapis terbentuk.
\textbf{2.} Di bawah 0.05 heterogenitas permukaan, dan di atas 0.30 kondensasi di
dalam pori, melanggar asumsi-asumsi BET.
\textbf{3.} \num{1.278e-3}, \num{2.410e-3}, \num{3.453e-3}, \num{4.621e-3}, \num{5.659e-3}, \num{6.813e-3} \unit{g.cm^{-3}}.
\textbf{4.} $v_{\mathrm m} = 1/(0.02205 + 0.000180) = \qty{45.0}{cm^3.g^{-1}}$.
\textbf{5.} $c = 1 + 0.02205/0.000180 = 124$: lapisan pertama terikat jauh lebih kuat
daripada cairannya.
\textbf{6.} Perpotongan itu sangat kecil dibandingkan dengan kemiringan kali $x$,
sehingga galat relatifnya besar; tetapi $v_{\mathrm m}$ bergantung pada jumlah kemiringan dan perpotongan,
yang didominasi oleh kemiringan.
\textbf{7.} $45.0/22\,414 = \qty{2.01e-3}{mol.g^{-1}}$.
\textbf{8.} $\num{2.01e-3} \times \num{6.022e23} \times \num{0.162e-18} = \qty{196}{m^2.g^{-1}}$.
\textbf{9.} Penyangga alumina: platina yang terpapar (Bagian III) luasnya hanya sekitar \qty{0.9}{m^2.g^{-1}}.
\textbf{10.} Pori-pori selebar molekul terisi penuh pada tekanan yang sangat
rendah, tidak lapis demi lapis, sehingga tidak ada \omterm{def:b3:surfaces-catalysis:coverage}{lapisan tunggal} yang dapat
dikenali.
\textbf{11.} $2 \times 0.205/22\,414 = \qty{1.83e-5}{mol.g^{-1}}$.
\textbf{12.} $0.0100/195.08 = \qty{5.13e-5}{mol.g^{-1}}$.
\textbf{13.} $D = 1.83/5.13 = 0.357$.
\textbf{14.} $a^3/4 = (0.3923)^3/4 = \qty{0.01509}{nm^3}$.
\textbf{15.} $\frac{\sqrt3}{4}(0.3923)^2 = \qty{0.0666}{nm^2}$, muka yang paling rapat; (100) dan (110) memberikan 0.0770 dan
\qty{0.1088}{nm^2}, dan rata-rata ketiga kerapatan situs memberikan \qty{0.0807}{nm^2}.
\textbf{16.} Sebuah bola memuat $\pi d^3/6v_{\mathrm{at}}$ atom, dan $\pi d^2/a_{\mathrm{at}}$ di antaranya berada di permukaan:
$D = 6v_{\mathrm{at}}/(a_{\mathrm{at}}d)$.
\textbf{17.} $d = 6 \times 0.01509/(0.0807 \times 0.357) = \qty{3.1}{nm}$.
\textbf{18.} Satu H per Pt permukaan; partikel bulat berukuran seragam (hasilnya
adalah rata-rata berbobot luas permukaan); seluruh platina tereduksi dan
dapat dijangkau; tidak ada hidrogen yang meluber ke penyangga.
\textbf{19.} $\num{2.0e-5}/\num{1.83e-5} = \qty{1.1}{s^{-1}}$.
\textbf{20.} $\num{2.0e-5}/0.0100 = \qty{2.0e-3}{mol.g^{-1}.s^{-1}}$ per gram platina.
\textbf{21.} $D = 1.12/10 = 0.112$ dan bukan 0.357: faktor 3.2.
\textbf{22.} Bahwa reaksinya peka struktur: situs-situs aktifnya (tepi, sudut,
muka tertentu) bukan fraksi permukaan yang tetap.
\textbf{23.} Frekuensi itu menghitung peristiwa per \omterm{def:b3:complex-kinetics:enzyme}{situs aktif}, sehingga
pengaruh banyaknya logam yang terpapar tersingkirkan.
\textbf{24.} \textbf{Diameter rata-rata partikel platina $\approx \qty{3.1}{nm}$.}
\end{solution}
